% id: 138-20
% book: 베이직쎈 확률과 통계 · 07 통계적 추정
% section: 유형 088 모평균과 표본평균의 차
% figure: none
% answer: 
% answer_source: 
% variant_level: 0 (원본 전사)
% 이 파일은 build.mjs 가 items.json 에서 만든다. 고칠 때는 items.json 을 고친다.
\smash{\raisebox{1.5mm}{\dmkichul{베이직쎈 확률과 통계 138쪽 20번}}}
정규분포 $\mathrm{N}(m{,}\mkern3mu\mathopen{}50^2)$을 따르는 모집단에서 크기가 $n$인 표본을 임의추출하여 구한 표본평균의 값이 $\overline{x}$이었다. 모평균 $m$을 신뢰도 $95\mkern3mu \%$로 추정할 때, $|m-\overline{x}|\le 5$가 되도록 하는 $n$의 최솟값을 구하시오. \dmcondfill{$\mathrm{P}(|Z|\le 2)=0.95$}{}
